\section{Conclusion}
\label{sec:conclusion}

We formulated sparse-PMU event attribution through competing physical
histories. Load/line predictors start from estimated past states and retain
the cause of the response after restoration. Active mass $M$ and its
increment $\Delta M$ expose the changing explanation rather than replacing
the physical dynamics.

The sealed continuous
evaluation gives \COneMixTimely{}/\COneMixTotal{} timely family/source
decisions with \COneMixExtras{} extra alarms for the mixture, and
\COneNomTimely{}/\COneNomTotal{} with \COneNomExtras{} for nominal.
The frozen ExtraTrees reference gives \COneETTimely{}/\COneETTotal{}
with \COneETExtras{} extras. Most mixture extras arise under out-of-bank
R/X, so the gain in timely attribution does not imply uniformly better
alarm control.

The earlier hackathon pipeline remains a broader nine-label,
feature-based system. Its 30-s-window validation is a separate result;
the causal ExtraTrees successor in this paper provides the direct
same-stream comparison. The physical bank's strongest supported gain is
timely load-source attribution, while learned features already diagnose
most lines in this catalog with fewer unmatched alarms.

The gain concerns a specified load/line catalog in one simulator. The
larger alarm burden under unmodeled R/X and the absence of an isolated
recovery-memory ablation limit any causal explanation of that gain. The
method does not yet establish all-family diagnosis, calibrated event
probabilities, field transfer or real-time deployment.
